% =========================================================================
% Chapter 3 - Fourier Neural Operators
% =========================================================================
\chapter{Fourier neural operators}
\label{ch:fno}

\section{Learning operators, not functions}

A conventional network learns a map between finite-dimensional vectors. Retrain
it for a different input size and nothing transfers. That is fatal here: the
grid $(N_1,N_2,N_3)$ of a plane-wave calculation is fixed by the cell and the
cutoff, so \emph{every material has a different one}.

A \emph{neural operator} instead learns a map between function spaces,
\begin{equation}
  \mathcal{G}: \mathcal{A} \longrightarrow \mathcal{U},
  \qquad
  a(\rr) \longmapsto u(\rr),
\end{equation}
and is then discretised only for evaluation. Its parameters describe the
operator, not a particular sampling of it, so one set of weights serves every
grid.

\section{Kernel integral operators}

The construction generalises the affine layer. Lift the input pointwise to a
higher-dimensional channel space, $v_0(\rr) = P\,a(\rr)$, apply $L$ layers of
the form
\begin{equation}
  v_{\ell+1}(\rr) = \sigma\!\left(
      W v_{\ell}(\rr) + b
      + \bigl(\mathcal{K}(v_\ell)\bigr)(\rr) \right),
  \label{eq:noperator}
\end{equation}
and project back pointwise, $u(\rr) = Q\,v_L(\rr)$. Here $W$ is a local
(channel-mixing) linear map and $\mathcal{K}$ is a \emph{kernel integral
operator}
\begin{equation}
  \bigl(\mathcal{K}v\bigr)(\rr) = \int_D \kappa(\rr,\rr')\,v(\rr')\,\mathrm{d}\rr' .
\end{equation}
The pointwise terms alone would give a network with no spatial coupling at all;
$\mathcal{K}$ is what makes it an operator.

\section{The Fourier layer}

Restricting the kernel to be translation invariant, $\kappa(\rr,\rr')
= \kappa(\rr-\rr')$, turns the integral into a convolution, which the
convolution theorem diagonalises:
\begin{equation}
  \bigl(\mathcal{K}v\bigr)(\rr)
    = \Fop^{-1}\!\left[ R(\GG)\cdot \Fop[v](\GG) \right](\rr),
  \label{eq:fourierlayer}
\end{equation}
where $R(\GG)$ is a learned complex multiplier. Rather than parameterising
$\kappa$ in real space, the Fourier neural operator parameterises $R$ directly
and \emph{truncates} it to the lowest $M_1{\times}M_2{\times}M_3$ modes.

Truncation does three things at once: it regularises (high frequencies are
where noise lives), it bounds the parameter count, and --- decisively --- it
makes the weight tensor independent of the grid. Nothing in $R$ refers to
$(N_1,N_2,N_3)$.

\subsection{Why this suits a crystal}

Three properties align the architecture with the physics of
Chapters~\ref{ch:ksdft}--\ref{ch:ofdft}.

\begin{description}
  \item[Periodicity is exact, not learned.] The FFT imposes Born--von Kármán
    boundary conditions, which a crystal unit cell already obeys. A padded
    convolutional network would have to spend capacity imitating that, and
    would still produce edge artefacts.
  \item[The coupling is global in one layer.] Multiplying in reciprocal space
    is convolving over the whole cell in real space. This directly addresses
    the non-locality of Section~\ref{sec:hk}: the Hartree kernel
    $4\pi e^2/G^2$ that the operator must partly represent is
    \emph{diagonal} in exactly the basis the layer works in.
  \item[Differential operators are free and exact.] On a plane-wave grid
    $\nabla \to i\GG$ and $\nabla^2 \to -G^2$ are exact for band-limited
    fields. Since the layer already computes FFTs, the physics-informed terms
    of Chapter~\ref{ch:pifno} cost almost nothing --- and carry no
    discretisation error that the loss would otherwise misattribute to the
    network.
\end{description}

\section{Discretisation invariance}
\label{sec:discinv}

For a real input the transform is taken with \texttt{rfftn}, halving the last
axis. The retained coefficients form four ``corners'' in the first two axes
--- combinations of low positive and low negative frequencies --- each with its
own weight block. The parameter count is
\begin{equation}
  4 \times C_{\mathrm{in}} \times C_{\mathrm{out}} \times M_1 M_2 M_3
  \quad\text{complex numbers,}
\end{equation}
independent of the grid.

Two further choices are required for weights to be \emph{meaningful} across
resolutions, not merely applicable.

\subsection{Normalisation}

The discrete transform must approximate the continuous Fourier series. With
coefficients
\begin{equation}
  c_{\GG} = \frac{1}{N}\sum_{n} u_n e^{-i\GG\cdot\rr_n},
\end{equation}
i.e. \texttt{norm="forward"} in both directions, $c_\GG$ converges to the
continuous series coefficient as $N\to\infty$ and its magnitude does not drift
with $N$. Under the default convention a $120^3$ field would enter a layer with
amplitudes $\sim125\times$ those of a $24^3$ field and shared weights would be
meaningless.

\begin{ptip}
Measured: a band-limited field sampled at $16^3$ and $32^3$ gives operator
outputs agreeing to a relative $1.5\times10^{-6}$ at the shared points.
\end{ptip}

\subsection{Mode selection}

Truncating at a fixed mode \emph{index} keeps a different physical wavevector
band for every cell size, since $|\GG_{\max}| = 2\pi M/L$. Truncating instead
at a fixed $G_{\max}$ --- retaining
\begin{equation}
  M_i = \left\lfloor \frac{G_{\max} L_i}{2\pi} \right\rfloor
\end{equation}
modes --- makes the operator act on the same band of physics in every material.
Poraqu\^e offers both; the physical selection is preferred when the data set
spans a wide range of cell sizes.

\section{Handling grids that differ between materials}

Three mechanisms together deliver true grid independence.

\begin{enumerate}
  \item \textbf{Dynamic mode truncation.} Weights are allocated for $M_i$
    modes; each forward pass uses
    $\min(M_i, \text{modes available on this grid})$ and leaves the rest
    untouched. A $24^3$ and a $120^3$ sample share the same parameters.
  \item \textbf{Resolution-invariant normalisation}, as above.
  \item \textbf{Shape-bucketed batching.} Samples of different spatial shape
    cannot be stacked. Rather than pad --- which wastes compute and injects
    fictitious vacuum into the FFT --- materials are grouped by
    $(N_1,N_2,N_3)$ and batches drawn within a group.
\end{enumerate}

Normalisation layers must also be shape-agnostic: \texttt{GroupNorm} is used
throughout, since batch statistics are meaningless when every sample has a
different spatial extent.

\section{Cell conditioning}

Two materials can share a grid shape and differ entirely in physical scale, so
the lattice must enter explicitly. Poraqu\^e supplies it twice: as three
fractional-coordinate channels appended to the input (bounded in $[0,1)$ for
every material), and through FiLM conditioning, in which a small network maps
rotation-invariant cell descriptors --- the three lattice lengths, the three
angle cosines, and $\Omega^{1/3}$ --- to per-channel scale and shift
parameters. Neither touches the spatial dimensions, so both are oblivious to
the grid.

\section{Architecture summary}

\begin{center}
\begin{tabular}{@{}ll@{}}
\toprule
Stage & Operation \\
\midrule
Input & $(B, 1, N_1, N_2, N_3)$, normalised field \\
Coordinates & append 3 fractional-coordinate channels \\
Lifting & $1\times1\times1$ convolution to width $C$ \\
Fourier layers ($\times L$) & Eq.~\eqref{eq:fourierlayer} + pointwise + GroupNorm + FiLM, residual \\
Projection & two $1\times1\times1$ convolutions \\
Output & $(B, 1, N_1, N_2, N_3)$ \\
\bottomrule
\end{tabular}
\end{center}

The residual connection around each Fourier layer stabilises deeper stacks; the
pointwise branch supplies the local, high-frequency content that mode
truncation removes from the spectral branch.

\section{The five parameters that size a run}
\label{sec:sizing}

Four settings fix the operator's shape and one fixes the grid it is trained on.
They are easy to confuse, because four of the five are ``some number of
channels'' and the fifth is the only one that is not a property of the model at
all. Written against the architecture above:

\begin{center}
\begin{tabular}{@{}llll@{}}
\toprule
Key & Symbol & Section & What it sets \\
\midrule
\opt{model.width}               & $C$   & \opt{model} & channels carried through the stack \\
\opt{model.modes}               & $M_i$ & \opt{model} & Fourier coefficients kept per axis \\
\opt{model.n\_layers}           & $L$   & \opt{model} & how many times the cycle runs \\
\opt{model.projection\_channels}& $C_P$ & \opt{model} & hidden width of the read-out \\
\opt{data.resolution}           & $N$   & \opt{data}  & longest grid axis after downsampling \\
\bottomrule
\end{tabular}
\end{center}

\subsection{The operator, with every index written out}
\label{sec:indexed}

Equations~\eqref{eq:noperator} and~\eqref{eq:fourierlayer} suppress the channel
indices, which is what makes the four architecture settings hard to tell apart:
they are all ``some number of channels'' until it is written down which sum
each one is the upper limit of. Written out in full, with $\rr$ a grid point
and $\kk$ a reciprocal-lattice point:

\paragraph{Lifting.}
\begin{equation}
  v^{(0)}_{o}(\rr) = \sum_{i=1}^{C_{\mathrm{in}}+3}
      W^{\mathrm{lift}}_{o i}\, x_{i}(\rr) + b^{\mathrm{lift}}_{o},
  \qquad o = 1,\dots,C
  \label{eq:idx-lift}
\end{equation}
The $+3$ is the appended fractional coordinates. No nonlinearity here: the
lift is affine.

\paragraph{Fourier layer $\ell = 1,\dots,L$.}
\begin{align}
  \hat v^{(\ell-1)}_{i}(\kk)
    &= \Fop\bigl[v^{(\ell-1)}_{i}\bigr](\kk)
    \label{eq:idx-fft} \\[2pt]
  \hat z^{(\ell)}_{o}(\kk)
    &= \begin{cases}
         \displaystyle\sum_{i=1}^{C}
           R^{(\ell,\,c(\kk))}_{i o\,k_1 k_2 k_3}\,\hat v^{(\ell-1)}_{i}(\kk)
           & \kk \ \text{retained},\\[10pt]
         0 & \text{otherwise},
       \end{cases}
    \label{eq:idx-spectral} \\[2pt]
  z^{(\ell)}_{o}(\rr)
    &= \Fop^{-1}\bigl[\hat z^{(\ell)}_{o}\bigr](\rr)
    \label{eq:idx-ifft} \\[2pt]
  y^{(\ell)}_{o}(\rr)
    &= \mathrm{GroupNorm}\Bigl(
         z^{(\ell)}_{o}(\rr)
         + \sum_{i=1}^{C} W^{(\ell)}_{o i}\, v^{(\ell-1)}_{i}(\rr)
         + b^{(\ell)}_{o} \Bigr)
    \label{eq:idx-norm} \\[2pt]
  y^{(\ell)}_{o}(\rr)
    &\leftarrow \bigl(1 + \gamma_{o}(\text{cell})\bigr)\, y^{(\ell)}_{o}(\rr)
                + \beta_{o}(\text{cell})
    \label{eq:idx-film} \\[2pt]
  v^{(\ell)}_{o}(\rr)
    &= v^{(\ell-1)}_{o}(\rr) + \sigma\bigl(y^{(\ell)}_{o}(\rr)\bigr)
    \label{eq:idx-residual}
\end{align}

\paragraph{Read-out.}
\begin{align}
  u_{p}(\rr)
    &= \sigma\Bigl(
         \sum_{o=1}^{C} P^{(1)}_{p o}\, v^{(L)}_{o}(\rr) + b^{(1)}_{p} \Bigr),
    \qquad p = 1,\dots,C_P
    \label{eq:idx-read1} \\
  \mathrm{out}_{q}(\rr)
    &= \sum_{p=1}^{C_P} P^{(2)}_{q p}\, u_{p}(\rr) + b^{(2)}_{q},
    \qquad q = 1,\dots,C_{\mathrm{out}}
    \label{eq:idx-read2}
\end{align}

\noindent
with the index ranges

\begin{center}
\begin{tabular}{@{}lll@{}}
\toprule
Index & Range & Set by \\
\midrule
$o,\,i$      & $1,\dots,C$                  & \opt{width} \\
$\ell$       & $1,\dots,L$                  & \opt{n\_layers} \\
$p$          & $1,\dots,C_P$                & \opt{projection\_channels} \\
$k_1,k_2,k_3$& $k_1<m_1,\ k_2<m_2,\ 0 \le k_3 < m_3$ & \opt{modes}, with \opt{resolution} \\
$c$          & $1,\dots,4$                  & the four $(\pm k_1, \pm k_2)$ corners \\
$q$          & $1,\dots,C_{\mathrm{out}}$   & the task \\
\bottomrule
\end{tabular}
\end{center}

\noindent
and the retained band, for a sample on an $(N_1,N_2,N_3)$ grid,
\begin{equation}
  m_1 = \min\!\left(M, \tfrac{N_1}{2}\right), \quad
  m_2 = \min\!\left(M, \tfrac{N_2}{2}\right), \quad
  m_3 = \min\!\left(M, \tfrac{N_3}{2}+1\right).
  \label{eq:retained}
\end{equation}
The third axis is the one \texttt{rfftn} halves, so it carries only
$k_3 \ge 0$; the first two carry both signs, which is what the four corners
$c$ enumerate.

\subsection{GroupNorm, FiLM, and where the nonlinearity goes}
\label{sec:blockorder}

Equations~\eqref{eq:idx-norm}--\eqref{eq:idx-residual} are three steps between
the inverse transform and the next layer's input, and their \emph{order} is
load-bearing. The canonical Fourier neural operator has none of them: Li
\emph{et al.} write $v_{\ell+1} = \sigma(W v_\ell + \mathcal{K}v_\ell)$, the
activation applied directly to the sum. Poraqu\^e inserts a normalisation and a
conditioning step in between, and closes with a residual.

\paragraph{GroupNorm --- because the batch is not a population.}
The spectral branch's output scale is not controlled by anything: it is a sum
over $m_1m_2m_3$ retained coefficients of a field whose spectrum varies from
material to material, so two samples in the same batch can leave
Eq.~\eqref{eq:idx-ifft} an order of magnitude apart. Something has to set the
scale the activation sees.

It cannot be \texttt{BatchNorm}. Batch statistics assume the samples in a batch
are draws from one population, and here they are not even the same
\emph{shape}: shape-bucketed batching (Section~\ref{sec:discinv}) means the
batch is whatever materials happen to share a grid, its size varies, and a
$24^3$ and a $120^3$ sample never meet. \texttt{GroupNorm} computes its
statistics \emph{within a sample}, over a group of channels and all of space,
so it is indifferent to both.

The group count is the largest divisor of $C$ not exceeding $8$ --- at $C=16$,
eight groups of two channels --- because \texttt{GroupNorm} requires
divisibility and \opt{width} is free.

\paragraph{FiLM --- the only place the lattice reaches the stack.}
Two materials can share a grid shape and differ entirely in physical scale, so
the cell must enter explicitly. Feature-wise linear modulation is one affine
map per channel, driven by the cell embedding:
\begin{equation}
  y^{(\ell)}_{o} \;\longleftarrow\;
  \bigl(1 + \gamma_{o}(\text{cell})\bigr)\, y^{(\ell)}_{o}
  + \beta_{o}(\text{cell}),
\end{equation}
which is Eq.~\eqref{eq:idx-film}. It touches no spatial index, so it is
oblivious to the grid --- and the projection producing $\gamma$ and $\beta$ is
\textbf{zero-initialised}, so at step zero $\gamma = \beta = 0$ and the whole
modulation is exactly the identity. Conditioning is something the operator
learns to use, not something imposed on it before it has learned anything.

\paragraph{Why FiLM comes after the normalisation and not before.}
Because a normalisation removes whatever a modulation just applied, to the
extent that the modulation is constant within a group. Applying the same
$(\gamma, \beta)$ on either side of \texttt{GroupNorm} and measuring the change
against no modulation at all:

\begin{center}
\begin{tabular}{@{}lrr@{}}
\toprule
Groups & Channels per group & Relative change \\
\midrule
\multicolumn{3}{@{}l}{\emph{modulation before the norm}} \\
$16$ & $1$  & $3\times10^{-6}$ \\
$8$  & $2$  & $0.32$ \\
$4$  & $4$  & $0.36$ \\
$1$  & $16$ & $0.39$ \\
\midrule
\multicolumn{3}{@{}l}{\emph{modulation after the norm --- what Poraqu\^e does}} \\
$8$  & $2$  & $\mathbf{1.16}$ \\
\bottomrule
\end{tabular}
\end{center}

With one channel per group the modulation is annihilated exactly: the
normalisation restores the mean and variance the affine map had just changed,
and only its \emph{sign} survives. At the shipped eight groups of two, roughly
a third of it survives, as within-group contrast. Placed after the norm it
acts in full. The ordering is not a stylistic preference --- it is the
difference between conditioning that works and conditioning that is mostly
undone one line later.

\paragraph{The residual comes last.}
$v^{(\ell)} = v^{(\ell-1)} + \sigma(\cdot)$ means each layer contributes a
correction rather than replacing the field, which is what makes a deeper stack
trainable, and it is why the identity is reachable: with FiLM at its
initialisation and $\sigma$ small, a fresh block is close to a no-op.

\begin{pnote}
\textbf{The activation between Fourier layers is \opt{silu} (or \opt{gelu}),
and it is already there.} $\sigma$ in Eq.~\eqref{eq:idx-residual} \emph{is}
\opt{model.activation}: one application per block, $L$ in total. What differs
from the canonical FNO is not whether the nonlinearity is present but where it
sits --- after the normalisation and the conditioning rather than immediately
after the spectral sum, so that it acts on a signal of controlled scale that
already knows which cell it belongs to. Applying it first would normalise the
activation's output instead of its input, which is the arrangement
normalisation exists to avoid.

The $\sigma$ in the read-out, Eq.~\eqref{eq:idx-read1}, is the same one, so the
stack contains $L+1$ nonlinearities and \opt{activation} governs all of them:
at the default \opt{silu}, every one of them is SiLU.
\end{pnote}

\subsection{Reading the cost off the indices}

Each setting appears in exactly one place, and that placement \emph{is} the
scaling:

\begin{description}
  \item[$C$ is the range of \emph{both} sums in Eq.~\eqref{eq:idx-spectral}.]
    It is the input index $i$ and the output index $o$ of the same
    contraction, so the weight tensor $R$ carries $C^2$ of them and the cost
    is \textbf{quadratic}. Nothing else in the architecture has this property.
  \item[$m_1 m_2 m_3$ are free indices on $R$, not summed.] Every retained
    coefficient gets its own $C \times C$ block, so the count is
    \textbf{cubic} in \opt{modes} --- $R$ has shape
    $(4, C, C, m_1, m_2, m_3)$, which is Eq.~\eqref{eq:idx-spectral} read as
    an array.
  \item[$L$ counts repetitions of the block, not terms in a sum.]
    Equations~\eqref{eq:idx-fft}--\eqref{eq:idx-residual} are applied $L$
    times with independent weights, so depth is \textbf{linear} --- and it
    composes the learned convolutions, which is how it buys range that
    \opt{modes} would charge cubically for.
  \item[$C_P$ appears only in Eqs.~\eqref{eq:idx-read1}--\eqref{eq:idx-read2},
    outside the loop.] It is the range of one sum, evaluated once, so it is
    \textbf{linear and paid once} --- $C C_P + C_P C_{\mathrm{out}}$ weights
    in total.
  \item[$N_1 N_2 N_3$ appears nowhere.] No index range above refers to the
    grid. That is Section~\ref{sec:discinv} restated: $\rr$ and $\kk$ are
    arguments, not indices, and \opt{resolution} enters only through
    Eq.~\eqref{eq:retained}, where it can \emph{cap} $m_i$ but never enlarge
    the weight tensor.
\end{description}

\begin{pnote}
$\sigma$ in Eq.~\eqref{eq:idx-residual} is \opt{model.activation} --- by
default \opt{silu}, optionally one of the per-channel learnable
Kolmogorov--Arnold variants --- and so is the read-out's nonlinearity in
Eq.~\eqref{eq:idx-read1}. One key, $L+1$ applications.

A per-channel variant is sized by the layer it sits in: $C$ learned functions
inside a Fourier block, $C_P$ in the read-out.

\opt{model.projection\_activation} overrides the read-out alone. Its default is
\opt{null}, meaning ``whatever \opt{activation} is'', which is the only thing a
reader should have to know. It exists because the read-out has not always
followed: it was a fixed $\mathrm{GELU}$ until 2026-08-28, so one key then
governed $L$ nonlinearities out of $L+1$ and the exception appeared in no
config. A checkpoint written before the change records no read-out
nonlinearity, and is reloaded with $\mathrm{GELU}$ on exactly that evidence ---
which is a statement about old files, not a default for new ones.
\end{pnote}

\subsection{What each one buys}

\begin{description}
  \item[\opt{width} --- how much information travels through.] The number of
    channels the field is carried in between the lifting layer and the
    read-out, and therefore how many distinct features of the density the
    operator can hold at once. It is the operator's representational capacity.
  \item[\opt{modes} --- how far it travels in space.] Mode $m$ on an axis of
    length $\ell_i$ is a wavelength $\ell_i/m$, so \opt{modes} fixes both the finest
    feature the learned convolution can resolve and the longest range over
    which it can correlate. It is a \emph{capacity limit}, not a requirement:
    on a grid too coarse to supply that many, fewer are used automatically,
    which is what lets one set of weights serve materials on different grids
    (Section~\ref{sec:discinv}).
  \item[\opt{n\_layers} --- how many times the cycle is applied.] Depth
    composes the learned convolutions. Two layers of $M$ modes reach further
    than one layer of $M$ modes, because the second acts on a field the first
    has already delocalised --- so depth buys effective range at \emph{linear}
    cost, which \opt{modes} does not.
  \item[\opt{projection\_channels} --- how the result is read out.] The hidden
    width of the two-convolution head that maps the $C$ channels back to
    the physical output. It appears once, at the end, never inside a Fourier
    layer, and is purely pointwise.
  \item[\opt{resolution} --- what the operator is shown.] Not an architecture
    parameter. It is the longest grid axis after the spectral downsampling of
    Section~\ref{sec:resample}, applied once when the cache is built, and it
    fixes the band of physics present in the data rather than the band the
    operator acts on.
\end{description}

\subsection{Where the parameters are}

The spectral tensor is the model. For the shipped configuration
($C=16$, $M=8$, $L=3$, $C_P=64$), counting a complex weight as two parameters
--- which is what \pyclass{FNO3d.n\_parameters} does, and therefore what the
training log prints:

\begin{center}
\begin{tabular}{@{}llrr@{}}
\toprule
Component & Where & Parameters & Share \\
\midrule
Spectral $R$    & Eq.~\eqref{eq:idx-spectral} & $3\,145\,728$ & $99.79\%$ \\
FiLM            & Eq.~\eqref{eq:idx-film}     & $3\,168$      & $0.10\%$ \\
Cell encoder    & $\gamma,\beta$ inputs       & $1\,312$      & $0.04\%$ \\
Read-out        & Eqs.~\eqref{eq:idx-read1}--\eqref{eq:idx-read2} & $1\,153$ & $0.04\%$ \\
Pointwise $W$   & Eq.~\eqref{eq:idx-norm}     & $816$         & $0.03\%$ \\
GroupNorm       & Eq.~\eqref{eq:idx-norm}     & $96$          & $0.00\%$ \\
Lifting         & Eq.~\eqref{eq:idx-lift}     & $80$          & $0.00\%$ \\
\midrule
Total           &                             & $3\,152\,353$ & \\
\bottomrule
\end{tabular}
\end{center}

One term in one equation is $99.79\%$ of the model. Cell conditioning and
normalisation together are $0.14\%$, and the read-out --- the whole of
\opt{projection\_channels} --- is $0.04\%$.

\noindent
Each row is its equation's index ranges multiplied out:
\begin{align}
  \text{spectral}  &: \; L \times 4 \times C^2 \times m_1 m_2 m_3
                      \quad\text{complex} \\
  \text{pointwise} &: \; L \times (C^2 + C) \\
  \text{lifting}   &: \; C\,(C_{\mathrm{in}} + 3) + C \\
  \text{read-out}  &: \; C C_P + C_P + C_P\,C_{\mathrm{out}} + C_{\mathrm{out}}
\end{align}
Only the first matters, and it grows as $C^2 M^3 L$ --- the $C^2$ from the two
summed channel indices, the $M^3$ from the three free mode indices, the $L$
from the repetition.

\subsection{Cost is not the parameter count}

Doubling each setting in turn, at the shipped configuration on a $32^3$ grid
(CPU, forward and backward, four threads):

\begin{center}
\begin{tabular}{@{}lrr@{}}
\toprule
Change & Parameters & Time \\
\midrule
\opt{width} $16 \to 32$               & $3.99\times$ & $2.89\times$ \\
\opt{modes} $8 \to 12$                & $3.37\times$ & $1.32\times$ \\
\opt{n\_layers} $3 \to 6$             & $2.00\times$ & $1.97\times$ \\
\opt{projection\_channels} $64 \to 128$ & $1.00\times$ & $1.31\times$ \\
\bottomrule
\end{tabular}
\end{center}

The two rows where the columns disagree are the ones worth remembering.

\opt{modes} is \emph{expensive in memory and cheap in time}: $(12/8)^3 = 3.375$
more weights, but the FFT is unchanged and only the contraction over retained
coefficients grows, which is not the dominant term at this size. It is the knob
that runs a model out of memory long before it runs it out of wall clock.

\opt{projection\_channels} is the reverse --- \emph{free in parameters and not
free in time}. It adds $1\,152$ weights, $0.07\%$ of the model, because it is
one pointwise map at the end; but that map is evaluated at every one of the
$N^3$ voxels, so its cost scales with the grid rather than with its own size.
Read the parameter count alone and it looks like nothing, which is why it is
often the best place to spend a budget that \opt{width} would otherwise consume
quadratically.

\subsection{Resolution}

Resolution is the only one of the five that changes the \emph{data}. At the
shipped architecture:

\begin{center}
\begin{tabular}{@{}lrr@{}}
\toprule
Grid & Voxels & Time \\
\midrule
$16^3$ & $0.12\times$ & $0.33\times$ \\
$24^3$ & $0.42\times$ & $0.68\times$ \\
$32^3$ & $1.00\times$ & $1.00\times$ \\
$48^3$ & $3.38\times$ & $2.36\times$ \\
$64^3$ & $8.00\times$ & $6.21\times$ \\
\bottomrule
\end{tabular}
\end{center}

Time approaches the $N^3\log N$ of the transforms as the grid grows, and falls
short of it below $32^3$ where fixed per-layer overhead dominates. Memory is
the harder limit: activations are $C N^3$ per layer and scale exactly, with no
overhead to hide behind.

\begin{pwarn}
\opt{resolution} and \opt{modes} are not independent. An $N^3$ grid supplies
$(N/2,\, N/2,\, N/2+1)$ modes --- the last axis is the halved one of
\texttt{rfftn} --- so at \opt{resolution: 32} any \opt{modes} above $16$
allocates weights that no sample can reach: allocated, saved into the
checkpoint, never trained. The shipped $M=8$ at $N=32$ uses half the available
band, which leaves room to raise one without the other.
\end{pwarn}

\begin{pnote}
None of the four architecture parameters refers to $N$. That is the point of
Section~\ref{sec:discinv}: the same weights are applied to a $24^3$ and a
$120^3$ sample, and \opt{resolution} can be changed without invalidating a
checkpoint's shape. What it does invalidate is the \emph{comparison} --- a
model trained at $32^3$ has seen a narrower band of physics than one trained at
$64^3$, whatever their reported errors say.
\end{pnote}

\subsection{Choosing them}

In one clause each: \opt{width} is how much information travels through;
\opt{modes} is how far it travels in space; \opt{n\_layers} composes the whole
cycle; and \opt{projection\_channels} is only how the result is read out.

A budget is usually best spent in that reverse order:

\begin{enumerate}
  \item \opt{projection\_channels}, first, because it is nearly free in
    parameters --- but judge it on the wall clock, not on the model size.
  \item \opt{n\_layers}, when the operator needs more \emph{range}: it buys
    that linearly, where \opt{modes} buys it cubically.
  \item \opt{width}, when it needs more \emph{capacity}, and pay the square.
  \item \opt{modes}, last and deliberately, after checking it against
    \opt{resolution}.
\end{enumerate}
